\subsection{トピックと参考文献の対応}

この表は、SWEBOKが示す「トピック対参考資料の行列」を、学習用に簡略化したものである。詳しく学ぶときは、各トピックに対応する文献を読む。

\small
\par\smallskip\noindent\refstepcounter{table}\label{tab:ch01-12}\textbf{表 \thetable: 第01章 ソフトウェア要求 表12}\par\smallskip
\begin{longtable}{@{}p{.38\linewidth}p{.55\linewidth}@{}}
\toprule
トピック & 主な参考資料 \\
\midrule
\endhead
1. ソフトウェア要求の基礎 & Wiegers and Beatty, \textit{Software Requirements}; Sommerville, \textit{Software Engineering} \\
1.1 要求の定義 & Wiegers and Beatty; Sommerville; ISO/IEC/IEEE 24765 \\
1.2 要求分類 & Wiegers and Beatty; Sommerville \\
1.3 製品要求とプロジェクト要求 & Wiegers and Beatty \\
1.4 機能要求 & Wiegers and Beatty; Sommerville \\
1.5 非機能要求 & Wiegers and Beatty; Sommerville; ISO/IEC 25010 \\
1.6 技術制約 & Tockey, \textit{How to Engineer Software}; Yourdon \\
1.7 サービス品質制約 & ISO/IEC 25010; Tockey, \textit{Return on Software} \\
1.8 分類の理由 & Wiegers and Beatty; Tockey; McMenamin and Palmer \\
1.9 システム要求とソフトウェア要求 & INCOSE Handbook; ISO/IEC/IEEE 29148 \\
1.10 派生要求 & Tockey; INCOSE Handbook \\
1.11 要求活動 & Wiegers and Beatty; Sommerville \\
2. 要求獲得 & Wiegers and Beatty; Sommerville; Robertson and Robertson; LaPlante and Kassab \\
2.1 要求源泉 & Wiegers and Beatty; Robertson and Robertson \\
2.2 要求獲得技法 & Wiegers and Beatty; Wood and Silver; Gottesdiener; Brown and Katz \\
3. 要求分析 & Wiegers and Beatty; Sommerville; Tockey \\
3.1 要求の性質 & Wiegers and Beatty; Robertson and Robertson \\
3.2 サービス品質要求の経済性 & Tockey, \textit{Return on Software} \\
3.3 形式的分析 & Sommerville; Wing; Software Engineering Models and Methods \\
3.4 要求衝突への対応 & Robertson and Robertson; Fisher and Ury; Weiss and Lai \\
4. 要求仕様化 & Wiegers and Beatty; Sommerville; ISO/IEC/IEEE 29148 \\
4.1 非構造化自然言語 & Wiegers and Beatty; ISO/IEC/IEEE 29148 \\
4.2 構造化自然言語 & Wiegers and Beatty; Cockburn; Robertson and Robertson \\
4.3 受け入れ基準に基づく仕様化 & Sommerville; Smart; Software Testing KA \\
4.4 モデルに基づく仕様化 & Wiegers and Beatty; Sommerville; Wing; Tockey \\
4.5 要求の追加属性 & Wiegers and Beatty; Gilb \\
4.6 増分的仕様化と包括的仕様化 & Wiegers and Beatty; Robertson and Robertson \\
5. 要求妥当性確認 & Wiegers and Beatty; Sommerville; LaPlante and Kassab \\
5.1 要求レビュー & Wiegers and Beatty; Software Quality KA \\
5.2 シミュレーションと実行 & Tockey; Model-Based Specification references \\
5.3 プロトタイピング & Wiegers and Beatty; Sommerville \\
6. 要求管理活動 & Wiegers and Beatty; Sommerville; McConnell \\
6.1 要求整理・削減 & McConnell, \textit{Rapid Development} \\
6.2 要求変更管理 & Wiegers and Beatty; Sommerville; Software Configuration Management KA \\
6.3 範囲調整 & McConnell; Software Engineering Management KA \\
7. 実務上の考慮 & Wiegers and Beatty; Sommerville; Tockey \\
7.1 反復性 & Sommerville; Robertson and Robertson \\
7.2 優先順位付け & Wiegers and Beatty \\
7.3 追跡可能性 & Wiegers and Beatty; Gotel and Finkelstein \\
7.4 安定性と変動性 & Sommerville; Tockey \\
7.5 要求測定 & Wiegers and Beatty; Software Engineering Process KA \\
7.6 プロセス品質と改善 & Wiegers and Beatty; Software Quality KA \\
8. ソフトウェア要求ツール & Wiegers and Beatty; Sommerville \\
8.1 要求管理ツール & Wiegers and Beatty \\
8.2 要求モデリングツール & Wiegers and Beatty; Sommerville; Formal Methods references \\
8.3 機能テストケース生成ツール & Software Testing KA; Tockey \\
\bottomrule
\end{longtable}
\normalsize
